\section{Conclusion}

This work introduces \textbf{TemporalBench}, a benchmark for systematically evaluating agent and LLM performance on time-series tasks that require temporal understanding, prediction, and contextual awareness.
By combining real-world time-series data with event-centric task construction and a four-tier task taxonomy (T1--T4), TemporalBench moves beyond traditional forecasting-only evaluations and enables structured assessment of temporal competence under diverse informational settings.

Our empirical results show that strong numerical forecasting performance does not reliably translate into correct decisions when tasks require contextual grounding or event-conditioned reasoning.
Performance varies substantially across task families and domains, indicating that temporal competence is not a single unified capability.
Improvements in forecasting accuracy alone are therefore insufficient to ensure robust behavior in settings that involve contextual interpretation or conditional prediction.

Taken together, these findings suggest that current agentic approaches to time-series problems exhibit fragmented strengths rather than holistic temporal understanding.
While some systems perform well on isolated aspects of temporal tasks, none consistently integrate temporal structure, contextual information, and decision requirements across all evaluated settings.

From a broader perspective, TemporalBench highlights the need for evaluation frameworks that disentangle complementary temporal competencies instead of collapsing them into a single end-to-end objective.
By evaluating agents across multiple task types, domains, and informational regimes, the benchmark provides a foundation for more principled analysis of agent capabilities and limitations in real-world time-series applications.
